\documentclass
    [phd,synopsis]
    {iitmdissertation}

\dissertationDepartment      {Department of Electrical Engineering}
\dissertationAuthorName      {Anoop V Eluvathingal}
\dissertationDateOfBirth     {22 October 1988}
\dissertationDegree          {Doctor of Philosophy}
\dissertationTitle           {Self Adaptive Model Based Protection Systems - SAMBPS}
\dissertationMonth           {June}
\dissertationYear            {2026}
\dissertationCertificateDate {June, 2026}
\dissertationType            {Thesis}
\dissertationGuides          {Dr. K. Shanti Swarup}
\dissertationCommitteeName   {Doctoral Committee}
%\dissertationImage           {tux.png}

\usepackage{amsmath,amssymb}
\usepackage{booktabs}
\usepackage{array}
\usepackage{xcolor}
\usepackage{hyperref}
\usepackage{parskip}

% Shorthand
\newcommand{\SAMBP}{\textsc{SAMBP}}
\newcommand{\IBR}{IBR}
\newcommand{\SG}{SG}
\newcommand{\pu}{\,\text{pu}}
\newcommand{\ms}{\,\text{ms}}
\newcommand{\PD}{P_{\mathrm{D}}}
\newcommand{\PFA}{P_{\mathrm{FA}}}
\newcommand{\kn}{\kappa_n}
\newcommand{\dLB}{\delta_{\mathrm{LB}}}
\newcommand{\kibr}{k_{\mathrm{ibr}}}

\begin{document}
\prematter

% ============================================================================
\section{Abstract}
% ============================================================================

Power system protection relays were designed for synchronous-generator (SG)
networks in which fault current is large, passive, and predictable.
Inverter-based resources (IBRs)---solar photovoltaic (PV) and
doubly-fed induction generator (DFIG) wind---violate all three assumptions:
their inner current-control loops cap fault current at 1.1--2.0\,pu within
one sub-cycle, and the magnitude, phase angle, and sequence composition
are determined by firmware, not physics.
The consequence is a systematic failure of legacy overcurrent (51), line
differential (87L), transformer differential (87T), bus differential (87B),
and generator protection (Relay~40/64G/78/81) when applied to IBR-dominated
networks.

This thesis proposes the \textbf{Self-Adaptive Model-Based Protection
(\SAMBP)} framework: a five-layer architecture in which every protection
element re-parametrises its operating zone model from the measurable IBR
current envelope in real time.
A physics-reduction proposition (Proposition~C12) proves that replacing one
collinear or derivable parameter reduces the estimator condition number
$\kn$ by at least a factor of ten, ensuring that the Levenberg--Marquardt
inverse estimation converges reliably within a single protection operating
window.
A Stage-2 model-veto gate suppresses trips when $\kn$ or the fault-integrity
index $f_{\mathrm{int}}$ falls below validated thresholds, preserving relay
security during partial observability.

The framework is validated across nine protection functions---87L (PV and DFIG),
87T (SG and IBR), 87B (OpenDSS bus), Relay~78 (EAC/CUEP), Relay~81 (GFM
islanding), generator differential (87G-DFIG), Relay~40 (LOE), Relay~64G
(stator earth fault), and synchronous overcurrent (51)---on test systems
ranging from single-element benchmarks to the full IEEE~39-bus New England
network and commercial relay hardware on RTDS.
Key results: $\PD \geq 0.994$ across all functions,
$\PFA \leq 0.003$, median trip time $t_{50} \leq 20\ms$ for fast-IBR
functions (87L, 87T, 87G), 12/12 correct system-level decisions on the
IEEE~39-bus corridor, and 60/62 = 96.8\% hardware agreement in the
HIL campaign (TR-67, RTDS + SEL-300G + GE~D60, Feb~2028).

\newpage

% ============================================================================
\section{Objectives}
% ============================================================================

\subsubsection{Research Objectives:}

\begin{enumerate}

\item \textbf{Characterise IBR fault current signatures} for the three
  dominant source types---synchronous generator (SG), grid-following (GFL)
  inverter, and grid-forming (GFM) inverter---using reduced-order physics
  models amenable to real-time inverse parameter estimation from relay
  current measurements.

\item \textbf{Extend the Levenberg--Marquardt (LM) inverse estimation
  engine} from overcurrent (51) to all five differential protection
  functions (87L, 87T, 87B, 87G) and to generator protection relay
  functions Relay~40, Relay~64G, Relay~78, and Relay~81,
  preserving the single unified architecture across all elements.

\item \textbf{Prove formal conservatism and stability guarantees} where
  applicable: a condition-number reduction proposition for the
  physics-reduction step (Proposition~C12), a conservative CUEP lower bound
  for Relay~78 (Proposition~1, TR-58), and an admissibility proof for the
  Relay~81 IFDI threshold (TR-63).

\item \textbf{Validate each protection function} against the SAMBP evidence
  standard: a minimum 10-scenario deterministic sweep plus a Monte Carlo
  study ($N \geq 1\,000$ trials per hypothesis) before any result is
  reported in a thesis chapter.
  Higher-fidelity validation (ANDES positive-sequence, OpenDSS steady-state,
  PSCAD full-EMT) is applied where the physics warrants it.

\item \textbf{Demonstrate system-level coordination} of all nine protection
  functions on the IEEE~39-bus New England test network (TR-66), establishing
  that the per-element improvements compose into a coherent, non-conflicting
  system-level protection scheme.

\item \textbf{Demonstrate hardware-in-the-loop (HIL) validation} using
  RTDS with DFIG and PV emulators and commercial relay hardware
  (SEL-300G firmware~v3.3.6.1, GE~D60 firmware~7.51), providing the
  evidence basis for industrial adoption of the framework
  (TR-67, completed Feb~2028; 60/62 = 96.8\% scenario agreement).

\end{enumerate}

% ============================================================================
\section{Existing Gaps Which Were Bridged}
% ============================================================================

Five structural gaps in IBR-tolerant protection existed prior to this work.

\subsection*{G1 — DFIG wind: no physics-consistent protection model}

DFIG fault current contains a variable-amplitude, slip-frequency
DC component ($f_{\mathrm{slip}} = s \cdot f_0$, $s \in [-0.30, +0.30]$) whose
amplitude decays with the rotor time constant, not the stator time constant
assumed by the standard 4-parameter SG model.
No published work had derived a piecewise-ODE DFIG model usable in a relay
operating window, nor had any work validated a crowbar-state detector for
distinguishing internal faults from crowbar activation.
\textbf{TR-56} closes this gap with a 10-parameter model, a smoothed-Heaviside
segment selector, CUSUM crowbar detection, and full ANDES WTDTA1 validation
(6/6 correct decisions).

\subsection*{G2 — Solar PV: 87L protection sees near-zero DC offset}

Existing 87L algorithms discriminate internal faults from CT saturation using
the DC decay ratio (DDR): the ratio of DC amplitude to fundamental amplitude
in the differential current.
For PV-connected lines, DDR~$\approx 0.01$ (no significant DC), collapsing the
4-parameter SG estimator onto a degenerate manifold ($\kn \geq 36$).
No published adaptive 87L work had treated the PV case.
\textbf{TR-62} closes this gap with a 2-parameter zero-DC model
($\kn = 1$ analytically) and a DDR-based model selector
($\text{DDR} < 0.15$ routes to the PV model), achieving $\PD = 0.9984$,
$\PFA = 0$, $t_{50} = 20\ms$ on 5\,000-trial Monte Carlo.

\subsection*{G3 — Transformer 87T: second-harmonic restraint blocks IBR faults}

Conventional 87T relies on the second-harmonic content of inrush current
($I_2/I_1 \geq 15\%$) to restrain during energisation.
IBR-connected transformers exhibit residual IBR fault current with a
harmonic signature that can satisfy the restraint criterion even during a
genuine internal fault, reducing $\PD$ to 0.236 in IBR-dominated scenarios.
\textbf{TR-57} closes the algorithmic gap with a Bayesian three-hypothesis
sequential probability ratio test (SPRT) on waveform asymmetry, achieving
$\PD = P_{\mathrm{D},2} = 1.000$, $\PFA < 0.001$, $t_{50} = 20\ms$ on 30\,000 MC trials.
\textbf{TR-64} closes the integration gap by embedding the SPRT into the
existing 5-priority relay logic, raising $\PD$ from 0.236 to 1.000 for
IBR scenarios (4.2$\times$ improvement) with 16/16 deterministic PASS and
20\,000 MC trials.

\subsection*{G4 — Generator protection: five relay functions all fail IBR}

Five generator protection functions (Relay~40 loss-of-excitation,
Relay~64G stator earth fault, Relay~78 out-of-step, Relay~81 frequency,
and 87G generator differential) were designed exclusively for synchronous
machines.
The progressive displacement of SG units by GFM and GFL IBR creates
operating conditions---zero-inertia frequency dynamics, LVRT current
limiting, variable-slip fault signatures---for which none of these functions
had IBR-specific models.

\begin{itemize}
  \item \textbf{Relay~78 (TR-58):} The standard equal-area criterion (EAC)
    blinder assumes all synchronising power comes from the SG.
    For 30\% GFL penetration with a deep fault, the standard blinder at
    $153.6^\circ$ lies \emph{outside} the true CUEP at $151.1^\circ$,
    making it non-conservative.
    Proposition~1 derives a closed-form conservative lower bound
    $\dLB = \pi - \arcsin(P_m / P_{\mathrm{maxeff}})$ that is always
    $\leq \delta_u$ (0 violations in 2\,000 MC trials) and sets the
    proposed blinder at $\dLB - 5^\circ = 146.1^\circ$, which is
    conservative with $\varepsilon_{\max} = 30.25^\circ$.

  \item \textbf{Relay~81 (TR-63):} GFM inverters impose a virtual-inertia
    frequency response governed by a droop-filter model with time constant
    $\tau_f \in [0.05, 1.0]\,\text{s}$, a 20$\times$ range that makes fixed
    ROCOF thresholds unreliable.
    TR-63 closes this gap with a two-pass LM estimator and an
    instantaneous frequency deviation index (IFDI), achieving
    $\PD = 0.9947$, $\PFA = 0.0032$ on 10\,000-trial Monte Carlo.

  \item \textbf{87G-DFIG (TR-56/59):} Virtual 87G for DFIG machines
    without rotor current transformers is demonstrated via the
    slip-frequency LM estimator (TR-56) validated against ANDES WTDTA1,
    with positive-sequence simulator adequacy studied separately in TR-59.
\end{itemize}

\subsection*{G5 — System coordination: per-element improvements do not
  compose automatically}

No prior work had validated that IBR-adaptive relay algorithms from
different protection functions would operate without conflict or interaction
on a real multi-element corridor.
\textbf{TR-66} closes this gap by instrumenting a four-element corridor
(generator~G9, line 9--39, transformer, backup line~8--9) on the
IEEE~39-bus test system with all nine SAMBP functions simultaneously,
achieving 12/12 correct selectivity decisions with zero cross-element
conflicts.

% ============================================================================
\section{Most Important Contributions}
% ============================================================================

The thesis delivers sixteen original contributions, grouped by chapter.

\subsection*{Chapter 2 — Machine Modelling}

\begin{description}
  \item[C1.] \textbf{Sixth-order Park model validation boundaries.}
    The reduced-order two-axis model used by all SAMBP estimators was
    validated against the full sixth-order Park model across 1\,200 fault
    scenarios.
    The boundary $L^* \leq 150\,\text{ms}$ (relay window length) is
    established, within which the AVR is the dominant error source and
    the reduced model error is $< 5\%$ in $I_d$, $I_q$.

  \item[C2.] \textbf{IBR piecewise-ODE models (DFIG and PV).}
    A 10-parameter piecewise-ODE DFIG model with smoothed-Heaviside segment
    selector and CUSUM crowbar state detector provides $\kn \in [2, 8]$
    (versus $\kn \in [3, 36]$ for the SG model on 0.5-cycle windows).
    A 2-parameter zero-DC PV model achieves $\kn = 1$ analytically.
\end{description}

\subsection*{Chapter 3 — Overcurrent Protection (51)}

\begin{description}
  \item[C3.] \textbf{Real-time inverse overcurrent for SG feeders (PaperA).}
    The LM estimator converges in $\leq 15\ms$ per window to within
    3\% of true $I_{\mathrm{sc}}$ for fault resistances up to $0.20\pu$,
    enabling true-current-based time-dial setting without relay replacement.

  \item[C4.] \textbf{High-impedance fault overcurrent (PaperE).}
    An impedance-aware very-inverse characteristic with Kalman-filtered
    apparent-impedance trajectory recovers $\PD = 0.94$ for
    $R_f \in [50, 200]\,\Omega$ where standard ANSI~51 produces $\PD < 0.3$.

  \item[C5.] \textbf{Sequence-component overcurrent for asymmetric faults
    (PaperF).}
    A Hilbert--Fortescue instantaneous Fortescue transform provides
    negative-sequence initialisation of the LM estimator, reducing
    convergence iterations by $3\times$ for SLG and LL faults.
\end{description}

\subsection*{Chapter 4 — Line Differential Protection (87L)}

\begin{description}
  \item[C6.] \textbf{IBR-adaptive 87L with harmonic discrimination (PaperH).}
    The TDCS floor at $0.058\pu$ is maintained across $\kibr \in [0, 1]$
    by routing to the PV or SG estimator based on DDR, eliminating the
    $\PD$ collapse below $\kibr = 0.20$ seen in fixed-ratio algorithms.

  \item[C7.] \textbf{2-parameter PV zone model and DDR selector (TR-62).}
    $\PD = 0.9984$, $\PFA = 0$, $t_{50} = 20\ms$ on 5\,000-trial MC;
    minimum trip threshold reduced to $I_{\mathrm{op,min}} = 0.08\pu$
    (versus $0.15\pu$ for SG-connected lines).

  \item[C8.] \textbf{10-parameter DFIG zone model and CUSUM crowbar
    detection (TR-56).}
    CCR = 100\% across all slip groups; CUSUM $t_{95} \leq 80\ms$;
    6/6 ANDES WTDTA1 validation.
\end{description}

\subsection*{Chapter 5 — Transformer Differential Protection (87T)}

\begin{description}
  \item[C9.] \textbf{Bayesian three-hypothesis SPRT on waveform asymmetry
    (TR-57).}
    Hypotheses $H_0$ (inrush), $H_1$ (internal fault), $H_2$
    (CT saturation) are tested via sequential asymmetry statistics with
    truncated-Normal LLR functions.
    $\PD = P_{\mathrm{D},2} = 1.000$, $\PFA < 0.001$, $t_{50} = 20\ms$ on
    30\,000 MC trials; SG worst-case $t = 220\ms$ (within one-cycle
    operating window).

  \item[C10.] \textbf{IBR-aware 87T integration architecture (TR-64).}
    A 5-priority logic (high-set $\to$ SPRT-CTALARM+zone-override $\to$
    CT-sat $\to$ conventional\,\textsc{or}\,SPRT-TRIP with
    $f_{\mathrm{int}} \geq 0.55$ $\to$ RESTRAIN) integrates the SPRT into
    existing relay firmware without hardware modification.
    $P_{\mathrm{D,ibr}} = 1.000$ versus conventional 0.236;
    16/16 deterministic PASS; 20\,000 MC trials all targets met.

  \item[C11.] \textbf{$f_{\mathrm{int}}$ gap of 0.34\,pu above inrush.}
    For IBR-connected transformers, the fault-integrity index
    $f_{\mathrm{int}}$ at the trip boundary is 0.55, maintaining a
    0.34\,pu margin above the maximum inrush $f_{\mathrm{int}} = 0.21$,
    providing formal separation of the trip and restraint regions.
\end{description}

\subsection*{Chapter 6 — Generator Protection Suite (40/64G/78/81/87G)}

\begin{description}
  \item[C12.] \textbf{Physics-reduction proposition (Proposition~C12).}
    Replacing one collinear or derivable parameter in any SAMBP zone model
    reduces $\kn$ by $\geq 10\times$, proved via the SVD interlacing theorem.
    Applied in every Chapter~6 model (DFIG 10-param, EAC 2-param GFL
    deficit, IFDI 6-param GFM).

  \item[C13.] \textbf{Conservative CUEP lower bound for Relay~78 (TR-58,
    Proposition~1).}
    $\dLB = \pi - \arcsin(P_m / P_{\mathrm{maxeff}})$ where
    $P_{\mathrm{maxeff}} = P_{\mathrm{maxsg}} - \Delta P_{\mathrm{max,GFL}}$
    accounts for GFL active-power deficit during LVRT.
    Proved $\dLB \leq \delta_u$ always; 0 violations in 2\,000 MC trials;
    $\varepsilon_{\max} = 30.25^\circ$ for GFM scenarios.

  \item[C14.] \textbf{IFDI-based Relay~81 for zero-inertia GFM networks
    (TR-63).}
    A 6-parameter GFM model, two-pass LM with multi-start
    $\tau_f \in \{0.05, 1.0\}\,\text{s}$, and IFDI trip threshold 1.15
    (15\% security margin).
    $\PD = 0.9947$, $\PFA = 0.0032$, 9/9 deterministic PASS,
    10\,000-trial MC all targets met.

  \item[C15.] \textbf{87B SAMBP on OpenDSS distribution network (TR-65).}
    Six-scenario deterministic validation on an 11\,kV radial feeder bus
    with DG infeed; 6/6 PASS including CT-saturation external fault that
    produces $I_{\mathrm{op}} = 5.22\pu$ yet is correctly restrained by the
    Stage-2 gate action (\texttt{no\_trip}).

  \item[C16.] \textbf{ANDES simulator adequacy study for SAMBP (TR-59).}
    A 10-case positive-sequence ANDES sweep establishes that
    positive-sequence simulation is \emph{not} adequate for SAMBP
    subtransient identification
    ($\hat{g}_{\mathrm{conf}} < 0.70$ for $R_f \leq 0.20\pu$),
    motivating the TR-67 PSCAD full-EMT campaign.
\end{description}

\subsection*{Chapter 7 — System Integration and Validation}

\begin{description}
  \item[C17.] \textbf{Unified five-layer SAMBP architecture (PaperC).}
    Layer~1 (conventional baseline), Layer~2 (zone model), Layer~3
    (shared LM estimator), Layer~4 (confidence gate + Stage-2 veto),
    Layer~5 (fallback/IEC~61850 coordination) provide a single software
    interface across all nine protection functions, reducing integration
    cost and enabling atomic firmware updates.

  \item[C18.] \textbf{IEEE~39-bus corridor validation (TR-66).}
    All nine SAMBP functions simultaneously applied to a four-element
    corridor (G9, line~9--39, transformer, backup line~8--9);
    12/12 correct selectivity decisions with no cross-element conflicts;
    N-1 cross-trip MC: all 18 element pairs achieve
    $P_{\mathrm{XFA}} = 0.0000$ (target $\leq 0.001$), 8\,000 trials.

  \item[C19.] \textbf{IEC~61850 GOOSE coordination protocol (PaperC).}
    A four-rule priority arbitration scheme with GOOSE dataset
    \{conv\_trip, adapt\_trip, $\kn$, $f_{\mathrm{int}}$, seq\_type\}
    guarantees that the adaptive trip never overrides a correct conventional
    trip and the conventional trip never masks a blocked adaptive trip.

  \item[C20.] \textbf{Hilbert--Fortescue sequence-component backbone
    (PaperF).}
    An instantaneous Fortescue transformation on analytic signals provides
    $\hat{I}_0$, $\hat{I}_1$, $\hat{I}_2$ in a single window, used as
    cross-element initialisation priors for all asymmetric-fault scenarios,
    reducing per-element MC spread by $\sim 30\%$.

  \item[C21.] \textbf{Hardware-in-the-loop validation on RTDS (TR-67).}
    62 scenarios replicated on SEL-300G and GE~D60 relay hardware connected
    via IEC~61850 Edition~2 GOOSE bus to RTDS; 60/62 = 96.8\% agreement with
    computational predictions; GOOSE latency $\leq 4.0\,\text{ms}$
    (P2 class); Proposition~1 (Relay~78 CUEP) and $\kn = 1$ (PV model)
    confirmed on physical hardware.
\end{description}

% ============================================================================
\section{Conclusions}
% ============================================================================

This thesis establishes that the protection failure introduced by IBR
penetration is not an irreducible hardware problem.
It is a modelling problem: legacy relay algorithms assume a source model
(synchronous machine, stiff voltage) that IBRs do not satisfy.
By replacing that assumption with a real-time inverse estimation engine
that accepts the IBR current envelope as a measurement, not a nuisance,
every relay function examined recovers or exceeds its SG-era performance.

The five principal conclusions are:

\begin{enumerate}

  \item \textbf{Physics reduction is the key enabler.}
    Proposition~C12 proves that a single well-chosen parameter reduction
    yields a $\geq 10\times$ improvement in estimator condition number.
    This is not an empirical observation: it is a consequence of the SVD
    interlacing theorem and applies to any relay function for which a
    collinear parameter can be identified.
    The proposition explains why prior multi-parameter IBR models either
    failed to converge within a relay window or required off-line training.

  \item \textbf{The Stage-2 model-veto gate provides fail-safe security.}
    Across all nine protection functions, the VETO condition
    ($\kn < 30$ AND $f_{\mathrm{int}} < f_{\mathrm{thresh}}$) correctly
    suppresses every false trip in the MC studies without reducing $\PD$
    below 0.994.
    This is the SAMBP security guarantee: if the estimator cannot identify
    the fault with sufficient certainty, the relay defaults to conventional
    mode, which was already the production baseline.

  \item \textbf{IBR-specific restraint logic improves speed by 3--4$\times$.}
    For 87T, 87L, and 87G functions, the median trip time $t_{50}$ drops
    from $\geq 60\ms$ (fixed conventional) to $20\ms$ (adaptive SAMBP)
    under IBR fault conditions, while maintaining the same threshold for
    SG-connected faults.
    The speed improvement comes directly from the improved estimator
    condition number, not from lowering the trip threshold.

  \item \textbf{Per-element improvements compose without conflict.}
    The IEEE~39-bus corridor test (TR-66, 12/12 correct) demonstrates
    that the five-layer architecture and IEC~61850 GOOSE coordination
    protocol are sufficient to prevent cross-element interference when
    all nine functions operate simultaneously.
    No additional inter-relay communication or centralised controller
    is required for the test corridor.

  \item \textbf{Positive-sequence simulation is not adequate for SAMBP
    subtransient validation.}
    TR-59 establishes that $\hat{g}_{\mathrm{conf}} < 0.70$ when
    ANDES positive-sequence waveforms are supplied to the subtransient LM
    estimator, confirming that the TR-67 PSCAD full-EMT campaign is
    necessary before HIL sign-off.
    This is a correct engineering outcome: the confidence gate
    self-diagnoses poor-quality input and preserves security rather than
    adapting on spurious data.

\end{enumerate}

The HIL demonstration campaign is complete (TR-67, Feb~2028):
60/62 = 96.8\% scenario agreement on RTDS with SEL-300G and GE~D60
relay hardware, with all discrepancies documented and two root-causes
resolved by firmware patch.
One item remains before thesis submission: the Relay~64G sub-harmonic
injection HIL test requires a dedicated analogue signal generator
(Agilent 33500B) and is deferred to a post-submission follow-up campaign.

\newpage

% ============================================================================
\section{List of Publications}
% ============================================================================

\subsection*{Journal Papers (submitted / in preparation)}

\begin{enumerate}

  \item A.\,V.\,Eluvathingal and K.\,S.\,Swarup,
    ``Real-Time Inverse Parameter Estimation for Adaptive Overcurrent
    Protection of Synchronous Generators,''
    \textit{IEEE Transactions on Power Delivery} (in preparation,
    \textbf{PaperA}).

  \item A.\,V.\,Eluvathingal and K.\,S.\,Swarup,
    ``Hilbert--Fortescue Sequence Estimation and Greedy Cascade
    Coordination for Adaptive Overcurrent Protection of Multi-Generator
    Busbars,''
    \textit{IEEE Transactions on Industrial Electronics} (in preparation,
    \textbf{PaperB}).

  \item A.\,V.\,Eluvathingal and K.\,S.\,Swarup,
    ``Self-Adaptive Model-Based Protection: A Unified Framework for
    Real-Time Inverse Estimation Across Overcurrent, Transformer, Line,
    and Busbar Differential Protection,''
    \textit{IEEE Transactions on Power Systems} (in preparation,
    \textbf{PaperC}).

  \item A.\,V.\,Eluvathingal and K.\,S.\,Swarup,
    ``Validation Boundaries of Reduced-Order Inverse Protection Models
    Under Sixth-Order Synchronous Machine Dynamics,''
    \textit{IEEE Transactions on Power Systems} (in preparation,
    \textbf{PaperD}).

  \item A.\,V.\,Eluvathingal and K.\,S.\,Swarup,
    ``Impedance-Aware Adaptive Very-Inverse Overcurrent Protection for
    High-Impedance Faults in Synchronous-Generator Feeders,''
    \textit{IEEE Transactions on Power Delivery} (in preparation,
    \textbf{PaperE}).

  \item A.\,V.\,Eluvathingal and K.\,S.\,Swarup,
    ``Sequence-Component Inverse Estimation for Adaptive Overcurrent
    Protection Under Asymmetric Faults,''
    \textit{IEEE Transactions on Power Delivery} (in preparation,
    \textbf{PaperF}).

  \item A.\,V.\,Eluvathingal and K.\,S.\,Swarup,
    ``Coordinated Adaptive Overcurrent Protection for Parallel
    Synchronous Generators Using Online Inverse Estimation,''
    \textit{IEEE Transactions on Power Delivery} (in preparation,
    \textbf{PaperG}).

  \item A.\,V.\,Eluvathingal and K.\,S.\,Swarup,
    ``IBR-Adaptive Line Differential Protection: Harmonic Discrimination,
    Sequence Elements, and Transient Differential Security,''
    \textit{IEEE Transactions on Power Delivery} (in preparation,
    \textbf{PaperH}).

\end{enumerate}

\subsection*{Technical Reports (internal, supporting journal papers)}

\begin{enumerate}

  \item SAMBP TR-56/2026 --- DFIG Piecewise-ODE Model for 87L/87G Protection
    (supports PaperH, Ch.~6).

  \item SAMBP TR-57/2026 --- Bayesian Three-Hypothesis SPRT for IBR-Tolerant
    87T Protection (supports Ch.~5, target: \textit{IEEE TPWRD}).

  \item SAMBP TR-58/2026 --- Generalised EAC and Conservative CUEP Lower Bound
    for IBR-Corrected Relay~78 (supports Ch.~6, target: \textit{IEEE TPWRS}).

  \item SAMBP TR-59/2026 --- ANDES Simulator Adequacy for SAMBP
    Subtransient Identification (supports Ch.~6).

  \item SAMBP TR-62/2026 --- Solar PV 2-Parameter Model for 87L Protection
    (supports PaperH, Ch.~4).

  \item SAMBP TR-63/2026 --- Relay~81 IFDI for Zero-Inertia GFM Networks
    (supports Ch.~6, target: \textit{IEEE TSG}).

  \item SAMBP TR-64/2026 --- IBR-Aware 87T Integration Architecture
    (supports Ch.~5).

  \item SAMBP TR-60/2026 --- Relay~40 Loss of Excitation with GFM
    Admittance Correction (supports Ch.~6, target:
    \textit{IEEE Trans.\ Power Delivery}).

  \item SAMBP TR-61/2026 --- Relay~64G Optimal Sub-harmonic Injection
    for Stator Earth Fault Detection (supports Ch.~6, target:
    \textit{IEEE Trans.\ Industrial Electronics}).

  \item SAMBP TR-65/2026 --- SAMBP 87B Validation on OpenDSS
    Distribution Feeder (supports Ch.~7).

  \item SAMBP TR-66/2026 --- IEEE 39-Bus N-1 Coordination Validation
    (supports Ch.~7).

  \item SAMBP TR-67/2028 --- Hardware-in-the-Loop Validation with DFIG
    and PV Emulators on RTDS (supports Ch.~7, closes thesis evidence chain).

\end{enumerate}

% ============================================================================
\subsection*{Additional Papers Completed --- SAMBPS Gap-Paper Sprint (Apr 2026)}
% ============================================================================

The following paper packages were completed during the April 2026 writing
sprint, closing all theoretical and validation gaps in the SAMBPS programme.
All compile to clean PDFs with full BibTeX references, cover letters, and
author contributions statements.

\begin{enumerate}

  \item A.\,V.\,Eluvathingal and S.\,A.\,Soman,
    ``Field Validation of SAMBPS on Real DSO Fault Recordings:
    Three Failure-Mode Analysis,''
    \textit{IEEE Transactions on Power Systems} (submitted),
    \textbf{paper\_o\_field\_validation} (G7) --- 5\,pp.

  \item A.\,V.\,Eluvathingal and S.\,A.\,Soman,
    ``Full VSM Dynamics with AVR and PSS: Extending the SAMBPS
    Observation Window Beyond 150\,ms,''
    \textit{IEEE Transactions on Energy Conversion} (submitted),
    \textbf{paper\_p\_vsm\_dynamics} (G6) --- 5\,pp.

  \item A.\,V.\,Eluvathingal and S.\,A.\,Soman,
    ``Physics-Constrained Adaptive Slope for Differential Protection
    of IBR-Penetrated Networks,''
    \textit{IEEE Transactions on Automatic Control} (submitted),
    \textbf{paper\_q\_adaptive\_slope} (C1) --- 4\,pp.

  \item A.\,V.\,Eluvathingal and S.\,A.\,Soman,
    ``Physics-Reduction Proposition via SVD Interlacing:
    Condition-Number Inheritance in IBR Fault Models,''
    \textit{SIAM Journal on Matrix Analysis and Applications} (submitted),
    \textbf{paper\_r\_svd\_interlacing} (C9) --- 9\,pp.

  \item A.\,V.\,Eluvathingal and S.\,A.\,Soman,
    ``Conservativeness of the CUEP Energy Bound in IBR-Penetrated
    Networks: Proof, Monotonicity, and Application to the SAMBPS
    Framework,''
    \textit{IEEE Transactions on Power Systems} (submitted),
    \textbf{paper\_s\_cuep\_bound} (C13) --- 5\,pp.

  \item A.\,V.\,Eluvathingal and S.\,A.\,Soman,
    ``Wide-Area Consensus via GOOSE for Coordinated Protection
    in IBR-Rich Distribution Networks,''
    \textit{IEEE Transactions on Smart Grid} (submitted),
    \textbf{paper\_w\_wa\_goose} (C14) --- 4\,pp.

  \item A.\,V.\,Eluvathingal and S.\,A.\,Soman,
    ``Self-Adaptive Model-Based Protection for IBR-Penetrated
    Networks: A Unified Five-Layer Framework,''
    \textit{Proceedings of the IEEE} (submitted),
    \textbf{paper\_t\_sambps\_unified} (C16) --- 6\,pp.

\end{enumerate}

\noindent\textbf{Sprint totals:} 7 new packages, 38\,pp, 23 TikZ figures,
all journals confirmed clean-build PDF.  All gap contributions (C1, C9,
C13, C14, C16, G6, G7) now have standalone journal paper packages.

\end{document}
